\documentclass[11pt]{article}
\usepackage[a4paper,margin=2.4cm]{geometry}
\usepackage{amsmath,amssymb,bm}
\usepackage{booktabs}
\usepackage{xcolor}
\usepackage[colorlinks=true,linkcolor=blue,urlcolor=blue]{hyperref}
\usepackage{enumitem}
\usepackage{siunitx}

\newcommand{\pd}[2]{\frac{\partial #1}{\partial #2}}
\newcommand{\ur}{u_r}
\newcommand{\uz}{u_z}
\newcommand{\rhalf}{r_{\mathrm{c}}}
\newcommand{\code}[1]{\texttt{\small #1}}
\definecolor{seedred}{RGB}{180,30,30}

\title{\bfseries Axisymmetric (RZ) Navier--Stokes in Cerisse:\\
Mathematical Analysis of the Near-Axis Discretization\\ and the Missing Viscous Geometric Terms}
\author{Cerisse SRP validation --- internal technical note}
\date{\today}

\begin{document}
\maketitle

\begin{center}
\fbox{\parbox{0.92\linewidth}{\textbf{Superseded regression status
(2026-07-19).}  This note records an earlier hypothesis and must not be used
as the current code assessment.  The cylindrical dilatation and hoop-stress
source are now present.  Direct RHS regressions show that the default WENO
pressure-flux plus $p_i/r_i$ pair is the only tested pressure treatment that
passes smooth cylindrical finite-volume tests; the pressure-split,
pressure-face and Simpson variants fail near the axis.  Current evidence and
remaining defects are recorded in
\code{exm/numerics/rz\_pressure\_operator},
\code{rz\_viscous\_axis}, and \code{rz\_planar\_shock}.}}
\end{center}

\begin{abstract}
\noindent
We derive the compressible Navier--Stokes equations in cylindrical
$(r,z)$ coordinates without swirl, cast them into the finite-volume
``metric divergence $+$ geometric source'' form that Cerisse uses, and
compare term-by-term against the implementation
(\code{src/tim/compute\_rhs.cpp}, \code{src/rhs/viscous.h}).  Two distinct
issues are identified and proven mathematically.
\textbf{(i) Euler level:} no physical term is missing --- the inviscid
geometric source $+p/r$ is present and is \emph{well balanced} for uniform
pressure.  However, the discrete radial-momentum balance is
\emph{not} well balanced \emph{across a shock}: the cell-centre pressure
used in the $+p/r$ source is inconsistent with the reconstructed face
pressures inside the flux, leaving a spurious radial-momentum source whose
amplitude is
$\bigl[p_{\mathrm{c}}-\tfrac12(p^{f}_{-}+p^{f}_{+})\bigr]/\rhalf$ and which
is amplified by $1/\rhalf=2/\Delta r$ in the on-axis cell. This seeds the
spurious near-axis velocity at on-axis shocks and scales like
$1/\Delta r$.
\textbf{(ii) Navier--Stokes level:} the cylindrical viscous geometric
terms are genuinely absent --- the dilatation omits $\ur/r$, the hoop
stress $\tau_{\theta\theta}$ is never formed, and the radial-momentum
source $-\tau_{\theta\theta}/r$ is not added (\code{rz\_geometric\_source}
is a no-op, and is additionally skipped for IBM/EB cases).  This does not
\emph{seed} the artifact (it is present in inviscid runs) but removes the
only physical sink that would damp it.  The defect lies in the NS/RZ
discretization, not in the IBM machinery.
\end{abstract}

\section{Governing equations in cylindrical coordinates}
Assume axisymmetry ($\partial/\partial\theta = 0$) and no swirl
($u_\theta=0$).  The velocity is $\bm u = \ur(r,z,t)\,\bm e_r +
\uz(r,z,t)\,\bm e_z$.  The compressible Navier--Stokes system reads
\begin{align}
\pd{\rho}{t} &+ \frac1r\pd{(r\rho\ur)}{r} + \pd{(\rho\uz)}{z} = 0,
\label{eq:cont}\\[2pt]
\pd{(\rho\ur)}{t} &+ \frac1r\pd{(r\rho\ur^2)}{r} + \pd{(\rho\ur\uz)}{z}
   + \pd{p}{r}
   = \underbrace{\frac1r\pd{(r\tau_{rr})}{r} + \pd{\tau_{rz}}{z}
     - \frac{\tau_{\theta\theta}}{r}}_{(\nabla\cdot\bm\tau)_r},
\label{eq:rmom}\\[2pt]
\pd{(\rho\uz)}{t} &+ \frac1r\pd{(r\rho\ur\uz)}{r} + \pd{(\rho\uz^2+p)}{z}
   = \frac1r\pd{(r\tau_{rz})}{r} + \pd{\tau_{zz}}{z},
\label{eq:zmom}\\[2pt]
\pd{E}{t} &+ \frac1r\pd{\bigl(r\ur(E+p)\bigr)}{r}
   + \pd{\bigl(\uz(E+p)\bigr)}{z}
   = \frac1r\pd{\bigl(r(\ur\tau_{rr}+\uz\tau_{rz}-q_r)\bigr)}{r}
   + \pd{\bigl(\ur\tau_{rz}+\uz\tau_{zz}-q_z\bigr)}{z}.
\label{eq:ene}
\end{align}
The Newtonian stress tensor (Stokes hypothesis, bulk viscosity $\xi$) has
the cylindrical components
\begin{equation}
\tau_{rr}=2\mu\,\pd{\ur}{r}+\lambda\Theta,\quad
\tau_{zz}=2\mu\,\pd{\uz}{z}+\lambda\Theta,\quad
\boxed{\;\tau_{\theta\theta}=2\mu\,\frac{\ur}{r}+\lambda\Theta\;},\quad
\tau_{rz}=\mu\Bigl(\pd{\ur}{z}+\pd{\uz}{r}\Bigr),
\label{eq:stress}
\end{equation}
with $\lambda=\xi-\tfrac23\mu$ and the \textbf{cylindrical dilatation}
\begin{equation}
\boxed{\;\Theta \;=\; \nabla\!\cdot\!\bm u
   \;=\; \pd{\ur}{r} + \frac{\ur}{r} + \pd{\uz}{z}\;}.
\label{eq:div}
\end{equation}

\paragraph{Where the geometric source terms come from.}
Because $\bm e_r$ rotates with $\theta$, the divergence operator carries
the metric factor $1/r$, and two terms in the radial momentum
equation~\eqref{eq:rmom} are \emph{not} expressible as a flux divergence
over the $(r,z)$ faces:
\begin{enumerate}[leftmargin=2.2em,itemsep=2pt]
\item the \textbf{inviscid pressure source} $+p/r$, obtained by writing the
      pressure inside the radial flux,
      $\tfrac1r\partial_r(rp)=\partial_r p + p/r$, so that
      \begin{equation}
      \pd{p}{r} \;=\; \frac1r\pd{(rp)}{r} \;-\; \frac{p}{r};
      \label{eq:psplit}
      \end{equation}
\item the \textbf{viscous hoop source} $-\tau_{\theta\theta}/r$ in
      $(\nabla\cdot\bm\tau)_r$.
\end{enumerate}
Continuity~\eqref{eq:cont}, axial momentum~\eqref{eq:zmom} and
energy~\eqref{eq:ene} are \emph{pure} flux divergences in cylindrical FV
form (the $\bm e_z$ direction and $u_\theta\!\equiv\!0$ produce no extra
source).  \emph{All} geometric source terms therefore live in the radial
momentum equation: one inviscid ($+p/r$) and one viscous
($-\tau_{\theta\theta}/r$).

\section{Finite-volume formulation used by Cerisse}
On a cell $[r_-,r_+]\times[z_-,z_+]$ with $r_\pm=r_0+(i,\,i{+}1)\Delta r$,
$\rhalf=\tfrac12(r_-+r_+)$, the cylindrical control volume is
$V=\pi(r_+^2-r_-^2)\Delta z$ and the radial face area is $A_r=2\pi r\Delta z$.
The metric-consistent radial divergence is
\begin{equation}
-\frac1r\pd{(rF_r)}{r}\;\Big|_{\text{cell}}
   \;\approx\;
   \frac{2\,(r_-F_- - r_+F_+)}{r_+^2-r_-^2}
   \;=\;
   \frac{r_-F_- - r_+F_+}{\rhalf\,\Delta r},
\label{eq:metricdiv}
\end{equation}
which is exactly \code{compute\_rhs.cpp:261}.  Cerisse builds the face
fluxes $F_\pm$ with the ordinary \emph{Cartesian} momentum flux
(WENO-Z5 characteristic Lax--Friedrichs); the RZ weighting is applied only
in \eqref{eq:metricdiv}.  The radial-momentum geometric source is then
added as
\begin{equation}
\Bigl(\pd{(\rho\ur)}{t}\Bigr)_{\!\text{source}} \mathrel{+}= \frac{p_{\mathrm c}}{\rhalf},
\qquad \frac{1}{\rhalf}=\frac{2}{\Delta r}\ \text{in the axis cell }(i{=}0),
\label{eq:psource}
\end{equation}
(\code{compute\_rhs.cpp:264--269}).  The viscous geometric hook
\code{rz\_geometric\_source} (\code{compute\_rhs.cpp:311}) is called only
when \emph{not} \code{AMREX\_USE\_GPIBM}/\code{CNS\_USE\_EB}, and its body is
empty (\code{src/rhs/viscous.h:587}, \code{viscousLES.h:628}).

\section{Euler level: correct, but not well balanced across a shock}
\label{sec:euler}
Write the discrete radial-momentum pressure contribution to the RHS,
using the face pressures $p^{f}_{\pm}$ that the flux actually carries
(from \eqref{eq:metricdiv}--\eqref{eq:psource}):
\begin{equation}
\mathcal R_p \;=\; \frac{r_-\,p^{f}_- - r_+\,p^{f}_+}{\rhalf\,\Delta r}
                 \;+\; \frac{p_{\mathrm c}}{\rhalf}.
\label{eq:Rp}
\end{equation}
The exact physical term is $-\partial p/\partial r \approx
-(p^{f}_+-p^{f}_-)/\Delta r$.

\paragraph{(a) Uniform pressure $\Rightarrow$ exactly balanced.}
With $p^{f}_-=p^{f}_+=p_{\mathrm c}=p$ and $r_--r_+=-\Delta r$,
\begin{equation}
\mathcal R_p = \frac{p\,(r_--r_+)}{\rhalf\Delta r} + \frac{p}{\rhalf}
             = -\frac{p}{\rhalf} + \frac{p}{\rhalf} = 0
             = -\pd{p}{r}\Big|_{\text{uniform}} .
\end{equation}
The $+p/r$ source exactly cancels the spurious $p/r$ generated by putting
$p$ in the radial flux (cf.\ \eqref{eq:psplit}).  Good.

\paragraph{(b) General pressure: the residual.}
Substituting $r_\pm=\rhalf\pm\tfrac12\Delta r$ into \eqref{eq:Rp} and
adding the exact $-\partial p/\partial r$, the discrete \emph{error}
(the spurious radial-momentum source) is, after cancellation,
\begin{equation}
\boxed{\;
\mathcal E \;\equiv\; \mathcal R_p - \Bigl(-\pd{p}{r}\Bigr)
   \;=\; \frac{\,p_{\mathrm c}-\tfrac12\bigl(p^{f}_-+p^{f}_+\bigr)\,}{\rhalf}\;}
\label{eq:residual}
\end{equation}
\emph{Derivation.}  Form $\mathcal R_p+(p^{f}_+-p^{f}_-)/\Delta r$ and
multiply by $\rhalf\Delta r$:
\[
r_-p^{f}_- - r_+p^{f}_+ + p_{\mathrm c}\Delta r + \rhalf\,(p^{f}_+-p^{f}_-)
= \tfrac12(r_--r_+)(p^{f}_-+p^{f}_+) + p_{\mathrm c}\Delta r
= \Delta r\Bigl[p_{\mathrm c}-\tfrac12(p^{f}_-+p^{f}_+)\Bigr],
\]
using $r_-p^{f}_- - r_+p^{f}_+ + \tfrac12(r_-+r_+)(p^{f}_+-p^{f}_-)
=\tfrac12(r_--r_+)(p^{f}_-+p^{f}_+)$ and $r_--r_+=-\Delta r$.  Dividing by
$\rhalf\Delta r$ gives \eqref{eq:residual}. \hfill$\square$

\paragraph{Consequence.}
\begin{itemize}[leftmargin=1.6em,itemsep=2pt]
\item \textbf{Smooth flow:} a (near-)central reconstruction gives
      $p_{\mathrm c}\approx\tfrac12(p^{f}_-+p^{f}_+)$, so
      $\mathcal E=\mathcal O(\Delta r^2)$ --- negligible.
\item \textbf{At a shock:} WENO reconstruction is upwind-biased and limited,
      so $p_{\mathrm c}\neq\tfrac12(p^{f}_-+p^{f}_+)$ and
      $\mathcal E=\mathcal O(\Delta p_{\text{shock}})/\rhalf$.  In the axis
      cell $\rhalf=\tfrac12\Delta r$, hence
      \begin{equation}
      \mathcal E_{\text{axis}} \;\sim\;
      \frac{2\,\bigl[p_{\mathrm c}-\overline{p^{f}}\bigr]}{\Delta r}
      \;\propto\; \frac{1}{\Delta r}.
      \label{eq:axisseed}
      \end{equation}
\end{itemize}
\textcolor{seedred}{This is the seed.}  A spurious radial-momentum source
of size $\mathcal O(\Delta p)/\Delta r$ appears wherever a strong shock
(the Mach disk, the bow shock) crosses the axis.  It violates the
symmetry condition $\ur(r{=}0)=0$, generates a near-axis radial velocity
and, through $\partial \ur/\partial z$, azimuthal vorticity.  No physical
PDE term is missing: \eqref{eq:axisseed} is a purely \emph{discrete}
well-balancedness defect of the singular coordinate, and it is
\emph{worse on finer grids}, $\propto 1/\Delta r$.

\section{Navier--Stokes level: the genuinely missing physics}
\label{sec:visc}
Cerisse evaluates the viscous stresses with the \emph{Cartesian}
dilatation (\code{viscous.h:256,472})
\begin{equation}
\Theta_{\text{code}} = \pd{\ur}{r}+\pd{\uz}{z}
   \;=\; \Theta - \frac{\ur}{r},
\end{equation}
i.e.\ the cylindrical $\ur/r$ term of \eqref{eq:div} is \textbf{dropped}.
Consequently every normal stress is in error by $\lambda\,\ur/r$:
\begin{equation}
\tau_{rr}^{\text{code}}=\tau_{rr}-\lambda\frac{\ur}{r},\quad
\tau_{zz}^{\text{code}}=\tau_{zz}-\lambda\frac{\ur}{r},\quad
\tau_{\theta\theta}^{\text{code}}=\text{(not formed)}.
\end{equation}
Moreover the hoop geometric source of \eqref{eq:rmom} is absent:
\begin{equation}
\boxed{\;\Bigl(\pd{(\rho\ur)}{t}\Bigr)_{\!\text{viscous-geom}}
  \mathrel{+}= -\frac{\tau_{\theta\theta}}{r}
  = -\frac{1}{r}\Bigl(2\mu\frac{\ur}{r}+\lambda\Theta\Bigr)
  \quad\text{is missing.}\;}
\label{eq:hoopmissing}
\end{equation}
Both omitted contributions carry a factor $1/r$ (directly, or through
$\ur/r$), so they are negligible in the bulk but become significant only
in the immediate neighbourhood of the axis.  The energy
equation~\eqref{eq:ene} needs \emph{no} separate geometric source --- the
viscous-work vector $\bm\tau\!\cdot\!\bm u$ has no $\theta$-component, so
$\nabla\!\cdot(\bm\tau\!\cdot\!\bm u)$ is captured by the metric divergence
--- but it inherits the wrong $\Theta_{\text{code}}$ in $\tau$.

\paragraph{Role.}  Equation~\eqref{eq:hoopmissing} is the \emph{only}
physical mechanism that diffuses near-axis azimuthal vorticity.  Its
absence does not create the artifact (the seed \eqref{eq:axisseed} is
inviscid), but it removes the sink that would otherwise damp it.

\section{Which defect causes the observed near-axis vortices?}
\label{sec:diag}
The two defects are independent and their roles are distinct:

\begin{table}[h]
\centering
\begin{tabular}{@{}lll@{}}
\toprule
 & \textbf{Euler seed \eqref{eq:axisseed}} & \textbf{Missing viscous hoop \eqref{eq:hoopmissing}}\\
\midrule
Nature & discrete well-balancedness defect & missing physical term \\
Active in & Euler \emph{and} NS runs & NS runs only \\
Effect & \emph{creates} spurious $\ur$ at on-axis shocks & \emph{fails to damp} it \\
Scaling & amplitude $\propto 1/\Delta r$ & --- \\
Fix & well-balanced pressure (route A) & implement \eqref{eq:div},\eqref{eq:hoopmissing} \\
\bottomrule
\end{tabular}
\end{table}

\noindent\textbf{Numerical evidence consistent with this analysis:}
\begin{itemize}[leftmargin=1.6em,itemsep=1pt]
\item Inviscid RZ jets (\code{no\_diffusive\_t}) already show a non-zero
      near-axis $\ur$ at the on-axis Mach disk $\Rightarrow$ the seed is
      \emph{inviscid}, ruling out \eqref{eq:hoopmissing} as the cause.
\item A Cartesian $M{=}6$ cylinder with the \emph{same} WENO-LF flux has a
      smooth, clean bow shock (HLLC, by contrast, carbuncles) $\Rightarrow$
      the flux family is not the culprit; the distinguishing factor is the
      RZ axis.
\item The artifact is stronger at finer resolution (prominent at $L4$,
      weak at $L2$), matching the $1/\Delta r$ scaling of
      \eqref{eq:axisseed}.
\end{itemize}
The bow-shock / Mach-disk \emph{stand-off distances} are unaffected,
because \eqref{eq:axisseed} perturbs the on-axis momentum balance, not the
global shock position.  The defect is in the NS/Euler RZ discretization;
the IBM ghost-point machinery acts in the solid-adjacent cells, away from
the open subsonic pocket where the vortices sit, and is not the driver.

\section{Implied corrections}
\begin{description}[leftmargin=1.4em,style=nextline]
\item[Route A --- well-balanced pressure (removes the seed).]
Discretize the radial pressure term so that \eqref{eq:residual} vanishes
by construction, i.e.\ remove $p$ from the radial-momentum flux and add a
\emph{cell-centred, consistent} gradient,
\begin{equation}
\Bigl(\pd{(\rho\ur)}{t}\Bigr)_{\!\text{pressure}}
  = -\,\frac{p^{f}_+-p^{f}_-}{\Delta r},
  \qquad p^{f}_\pm \ \text{the same face pressures used in the flux,}
\end{equation}
which is well balanced for arbitrary $p$ and carries no $1/r$
amplification.
\item[Route C --- complete the RZ viscous operator (restores physics).]
Use the cylindrical dilatation \eqref{eq:div} in all stresses, form
$\tau_{\theta\theta}$, add the hoop source \eqref{eq:hoopmissing} to the
radial momentum, and remove the \code{AMREX\_USE\_GPIBM}/\code{CNS\_USE\_EB}
gate at \code{compute\_rhs.cpp:311}.  This provides the physical damping of
near-axis vorticity and makes RZ runs true Navier--Stokes.
\end{description}

\paragraph{Summary.}
At the Euler level nothing physical is missing; the spurious near-axis
velocity is a discrete \emph{well-balancedness} failure of the $+p/r$
geometric source across shocks, eq.~\eqref{eq:residual}, amplified by
$2/\Delta r$ on the axis.  At the Navier--Stokes level the cylindrical
viscous terms ($\ur/r$ in $\Theta$, $\tau_{\theta\theta}$, and
$-\tau_{\theta\theta}/r$) are genuinely absent, eq.~\eqref{eq:hoopmissing},
removing the physical sink for the seeded perturbation.  Both are
properties of the NS/RZ discretization, not of the immersed boundary.

\end{document}
